\documentclass[border=5pt]{standalone}
\usepackage{circuitikz}
\usepackage{siunitx}

\begin{document}
\begin{circuitikz}[american, thick]
    % Fuente de voltaje
    \draw (0,0) to[V, l=\SI{10}{V}] (0,4);
    
    % Líneas principales
    \draw (0,4) -- (2,4);
    \draw (0,0) -- (8,0);
    
    % Amperímetro total (opcional, para medir la corriente total)
    \draw (2,4) to[rmeter, t=A, l=$I_{total}$] (4,4);
    \draw (4,4) -- (8,4);
    
    % Rama 1: 1k ohm
    \draw (4,4) to[R, l=$R_1$, a=\SI{1}{k\ohm}] (4,2);
    \draw (4,2) to[short, i=$I_1$] (4,0);
    
    % Rama 2: 330 ohm
    \draw (6,4) to[R, l=$R_2$, a=\SI{330}{\ohm}] (6,2);
    \draw (6,2) to[short, i=$I_2$] (6,0);
    
    % Rama 3: 220 ohm
    \draw (8,4) to[R, l=$R_3$, a=\SI{220}{\ohm}] (8,2);
    \draw (8,2) to[short, i=$I_3$] (8,0);
    
\end{circuitikz}
\end{document}
